\documentclass[a4paper]{article}
\usepackage{amsfonts}
\usepackage{amsmath}
\usepackage{amssymb}
\usepackage{graphicx}     

\begin{document}
  
\noindent \large Auteur: Zaky Souai \\
\noindent \large Studierichting:  Informatica\\
\noindent \large Oefeningenreeks 1D nr. 21.14\\

\bigskip

\normalsize

$z^4 - 5z^3 + 13z^2 - 19z + 10 = 0$\\

We proberen 1 uit:\\

$1 - 5 + 13 - 19 + 10 = 0 $\\

$\Rightarrow 1 $ is een nulpunt. Door Horner toe te passen, krijgen we: \\

$\Rightarrow (z^3 - 4z^2 + 9z - 10)(z-1)$\\

We proberen 2 uit:\\

$8 - 16 + 18 - 10 = 0$\\

$\Rightarrow 2$ is een nulpunt van $z^3 - 4z^2 + 9z - 10$. Door Horner weer toe te passen, krijgen we: \\

$(z^2 - 2z + 5)(z-1)(z-2)$\\

We zoeken nu nog de nulpunten van $z^2 - 2z + 5$:\\

$D = 4 - 20 = -16$\\

$z = \dfrac{2\pm{4i}}{2} = 1 \pm{2i}$\\

Antwoord: $1,2,1-2i,1+2i$\\


\begin{thebibliography}{999}
%\bibitem{a} Referentie
\end{thebibliography}
\end{document}